\documentclass[11pt]{article}
\usepackage{amsmath,amssymb,amsthm}
\usepackage{booktabs}
\usepackage{enumitem}
\usepackage{hyperref}
\setlength{\textwidth}{6.5in}\setlength{\oddsidemargin}{0in}\setlength{\evensidemargin}{0in}
\setlength{\textheight}{9.4in}\setlength{\topmargin}{-0.7in}
\setlist{nosep,leftmargin=1.6em}
\makeatletter
\renewcommand\section{\@startsection{section}{1}{\z@}{-2.2ex plus -.5ex minus -.2ex}{1.1ex plus .2ex}{\normalfont\large\bfseries}}
\renewcommand\paragraph{\@startsection{paragraph}{4}{\z@}{1.2ex plus .5ex minus .2ex}{-1em}{\normalfont\normalsize\bfseries}}
\makeatother
\AtBeginDocument{\setlength{\abovedisplayskip}{4pt plus 2pt}\setlength{\belowdisplayskip}{4pt plus 2pt}}
\newtheorem{theorem}{Theorem}
\newtheorem{lemma}{Lemma}
\newtheorem{proposition}{Proposition}
\newtheorem{fact}{Fact}
\newtheorem{conjecture}{Conjecture}
\theoremstyle{definition}
\newtheorem{definition}{Definition}
\theoremstyle{remark}
\newtheorem{remark}{Remark}
\newcommand{\cols}{\mathrm{cols}}
\newcommand{\ports}{\mathrm{ports}}
\newcommand{\inp}{\mathrm{in}}
\newcommand{\out}{\mathrm{out}}
\newcommand{\rest}{\mathrm{rest}}
\newcommand{\env}{\mathrm{env}}
\newcommand{\Al}{\mathop{\Box}}
\newcommand{\Ev}{\mathop{\Diamond}}
\newcommand{\eps}{\varepsilon}
\newcommand{\R}{\mathbb{R}}
\newcommand{\Q}{\mathbb{Q}}
\title{Decomposing a global PolySTL contract over an architecture graph}
\author{}
\date{18 October 2026}
\begin{document}
\maketitle

\section{Problem}
The problem, verbatim:
\begin{quote}\small
We have a system which is a graph of components where each component has labeled/typed
inputs and outputs (``typed'' means the variables are distinct and each is assigned to a known
component; there are no numeric range bounds). We have a global query contract written in PolySTL (polytopic signal
temporal logic: STL whose predicates are convex polytopes over the state variables,
mu(x) := A x $<=$ b, where x is the vector of ALL system variables).

We need to ALGEBRAICALLY DECOMPOSE the global contract into obligations for the components.
The delegate contract handed to a component may ONLY reference the variables that component
has available, i.e. its own inputs and outputs.

Only some components will receive a (non-trivial) delegated contract: not every global
contract needs every component.

All that is available is: (1) the architectural graph (components, typed/labelled ports,
who reads/writes which variable), and (2) the top-level query contract. No component models.

How it is done is open. It must be done.
\end{quote}
We work in discrete time on a bounded horizon, with free trajectories. Section~\ref{sec:repro} gives the one command that recomputes every formula and number below.

\section{Definitions}\label{sec:defs}
\paragraph{Architecture.}
$V$ is a finite set of real-valued variables. $C$ is a finite set of components.
Each variable has exactly one owner: a component or the environment.
For a component $i$, $\out(i)$ is the set of variables it owns and $\inp(i)\subseteq V\setminus\out(i)$ is the set it reads.
Put $\ports(i):=\inp(i)\cup\out(i)$. Each component has a delay $d_i\in\{0,1\}$.
There is an edge $i\to j$ when $\out(i)\cap\inp(j)\neq\emptyset$.
No numeric range is attached to any variable.

\paragraph{Traces.}
$H\in\mathbb{N}$ is the horizon and $T:=\{0,\dots,H\}$.
A column is a pair $(v,t)\in V\times T$, written $v(t)$. For $W\subseteq V$ let $\cols(W):=W\times T$.
A trace is a map $x:V\times T\to\R$.
A formula is a set of traces. It is \emph{over} a column set $K$ if membership depends only on $x|_K$.
Boolean connectives are set operations. $\varphi\models\psi$ means $\varphi\subseteq\psi$.

\paragraph{PolySTL.}
Formulas are $\varphi::=a^\top x\le b\mid a^\top x<b\mid\neg\varphi\mid\varphi\wedge\varphi\mid\varphi\vee\varphi\mid
\Al_{[\alpha,\beta]}\varphi\mid\Ev_{[\alpha,\beta]}\varphi$, with $a\in\Q^V$, $b\in\Q$ and integers $0\le\alpha\le\beta$.
The atom $a^\top x\le b$ holds at time $t$ iff $\sum_v a_v\,x(v,t)\le b$.
$\Al_{[\alpha,\beta]}\varphi$ holds at $t$ iff $\varphi$ holds at every $s\in[t+\alpha,t+\beta]$; $\Ev$ asks for some $s$.
We use only formulas whose evaluation at time 0 needs times $\le H$. A trace satisfies $\varphi$ iff $\varphi$ holds at time 0.
A \emph{time-stamped atom} is a constraint $\sum_v a_v\,x(v,t)\le b$ (or $<b$) at one time $t$. We write it $(a^\top x)(t)\le b$, for example $m(1)\ge 2$.
When $H=0$ we drop the time.
$|e|\le c$ abbreviates $e\le c\wedge -e\le c$, and $|e|>c$ abbreviates $e>c\vee -e>c$.

\begin{lemma}[unrolling]\label{lem:unroll}
Every PolySTL formula is a Boolean combination of time-stamped atoms. Conversely the time-stamped atom $(a^\top x)(t)\le b$ is the PolySTL formula $\Ev_{[t,t]}(a^\top x\le b)$.
\end{lemma}
\begin{proof}
By induction on the formula. An atom evaluated at $t$ is a time-stamped atom. Boolean connectives are kept.
$\Al_{[\alpha,\beta]}$ and $\Ev_{[\alpha,\beta]}$ at $t$ become a finite conjunction and disjunction over $s\in[t+\alpha,t+\beta]$.
The converse is the semantics of $\Ev_{[t,t]}$.
\end{proof}
So a \emph{PolySTL formula over $K$} is the same as a Boolean combination of time-stamped atoms whose columns lie in $K$.
A cross-time atom, which relates columns at two different times, is not of this kind.

\paragraph{Contracts.}
A contract for component $i$ is a pair $(A_i,G_i)$ of PolySTL formulas over $\cols(\ports(i))$.
Its implementation set is $S_i:=\neg A_i\vee G_i$.
The query is a pair $P=(A_P,G_P)$ of PolySTL formulas over $\cols(V)$.
A delegation is a family $(A_i,G_i)_{i\in C}$.
\begin{definition}
A delegation is \emph{valid} if
(C1) $A_P\wedge\bigwedge_i S_i\models G_P$ and (C2) $A_P\wedge\bigwedge_i S_i\models A_j$ for every $j$.
Component $i$ is \emph{receptive} if for every value $u$ of $\cols(\inp(i))$ some value $w$ of $\cols(\out(i))$ gives $(u,w)\in S_i$.
Let $u(t)$ and $w(t)$ be the input and output columns of $i$ at time $t$.
Component $i$ is \emph{causal} if the sentence $\mathcal{Q}_0\cdots\mathcal{Q}_H.\,S_i$ holds, where $\mathcal{Q}_t:=\forall u(t)\,\exists w(t)$ if $d_i=0$ and
$\mathcal{Q}_t:=\exists w(t)\,\forall u(t)$ if $d_i=1$.
By Skolemisation this says that each $w(t)$ can be computed from inputs at times $\le t-d_i$ and from earlier outputs.
A delegation is \emph{acceptable} if it is valid and every component is receptive and causal.
\end{definition}

We use the assume--guarantee contracts of Benveniste et al.~\cite{contracts}, lifted to traces over all columns.
The saturated form of $(A,G)$ is $(A,\neg A\vee G)$.
Two saturated contracts compose as $(A_1,S_1)\otimes(A_2,S_2):=((A_1\wedge A_2)\vee\neg(S_1\wedge S_2),\,S_1\wedge S_2)$.
Refinement is $(A,S)\preceq(A',S')$ iff $A'\subseteq A$ and $S\subseteq S'$.
\begin{theorem}\label{thm:refine}
A delegation is valid iff $\bigotimes_{i\in C}(A_i,S_i)\preceq(A_P,\neg A_P\vee G_P)$.
\end{theorem}
\begin{proof}
First, $\bigotimes_{i\in C}(A_i,S_i)=(\bigwedge_iA_i\vee\neg\bigwedge_iS_i,\ \bigwedge_iS_i)$.
This holds for one factor. If it holds for a family $\mathcal{F}$ with assumption $A_{\mathcal{F}}\vee\neg S_{\mathcal{F}}$ and guarantee $S_{\mathcal{F}}$, add a factor $(A',S')$. The guarantee becomes $S_{\mathcal{F}}\wedge S'$.
The assumption becomes $((A_{\mathcal{F}}\vee\neg S_{\mathcal{F}})\wedge A')\vee\neg(S_{\mathcal{F}}\wedge S')=(A_{\mathcal{F}}\wedge A')\vee(\neg S_{\mathcal{F}}\wedge A')\vee\neg(S_{\mathcal{F}}\wedge S')$, and the middle term is contained in the last.
Now compare with $(A_P,\neg A_P\vee G_P)$.
The guarantee condition $\bigwedge_iS_i\subseteq\neg A_P\vee G_P$ is (C1).
The assumption condition $A_P\subseteq\bigwedge_iA_i\vee\neg\bigwedge_iS_i$ says $A_P\wedge\bigwedge_iS_i\subseteq\bigwedge_iA_i$, which is (C2).
\end{proof}

\paragraph{Projections.}
For a formula $\varphi$ and a column set $K$, the universal and the existential projection are
\[\Pi_K\varphi:=\{x:\ \forall y\,(y|_K=x|_K\to y\in\varphi)\},\qquad
\pi_K\varphi:=\{x:\ \exists y\,(y|_K=x|_K\wedge y\in\varphi)\}.\]
Both are over $K$. Also $\Pi_K\varphi\subseteq\varphi\subseteq\pi_K\varphi$ and $\Pi_K\varphi=\neg\pi_K\neg\varphi$.
We compute both by linear quantifier elimination with the Z3 tactic \texttt{qe2}~\cite{z3,qe2}.
The script checks every result against its quantified input by a Z3 validity query.
\begin{fact}[PolySTL closure]\label{fact:closure}
If $\varphi$ is a Boolean combination of time-stamped atoms, so are $\Pi_K\varphi$ and $\pi_K\varphi$.
Each atom of the result lies at the time of some atom of $\varphi$.
\end{fact}
\begin{proof}
Since $\Pi_K=\neg\pi_K\neg$ and negation keeps the class, it suffices to treat $\pi_K$ and to eliminate one column $c=v(t)\notin K$.
Write $\varphi$ in disjunctive normal form. A negated atom is an atom: $\neg(a^\top x\le b)$ is $-a^\top x<-b$.
The quantifier $\exists c$ distributes over $\vee$, so take one conjunction of atoms.
Atoms without $c$ are kept. An atom with $c$ lies at time $t$, because all columns of a time-stamped atom share one time.
Solved for $c$, it reads $c\le u$ or $c\ge\ell$ (strict or not), with $u,\ell$ affine in the other columns at time $t$.
For fixed values of those columns, the admissible $c$ form an interval of the reals.
It is nonempty iff $\ell\le u$ for every lower bound $\ell$ and upper bound $u$, strict when either bound is strict.
So $\exists c$ of the conjunction is the kept atoms together with the atoms $\ell\le u$ (Fourier--Motzkin).
Each new atom combines two atoms that share $c$, so it lies at time $t$.
Repeating over the eliminated columns proves the claim. The result does not depend on eliminated columns, so it is over $K$.
\end{proof}

\section{Why the projections alone cannot choose}\label{sec:choose}
The relay is the running example. The environment owns $r$. Radio reads $r$ and owns $m$. Display reads $m$ and owns $d$.
Both delays are 0 and $H=0$, so columns are variables. The query is $A_P=|r|\le1$, $G_P=|d-r|\le\tfrac12$.
\begin{lemma}\label{lem:empty}
$\Pi_{\{m,d\}}(A_P\to G_P)=\mathit{false}$ and $\Pi_{\{r,m\}}(A_P\to G_P)=\neg A_P$.
\end{lemma}
\begin{proof}
A value $d$ in the first projection needs $|d-r|\le\frac12$ for $r=1$ and for $r=-1$, so $d\ge\frac12$ and $d\le-\frac12$.
For the second: if $|r|\le1$, the value $d=r+1$ violates $G_P$; if $|r|>1$ the implication holds.
The script prints \texttt{false} and \texttt{|r| > 1}.
\end{proof}
So projecting the query onto one component's ports gives an obligation that fails on $A_P$.
The edge $m$ must carry information about $r$. What it carries, and how precisely, is not fixed by the projections.
Order delegations by $D\sqsubseteq D'$ iff $S_i\subseteq S'_i$ for every $i$.
\begin{proposition}\label{prop:nogreatest}
The acceptable delegations of the relay have no $\sqsubseteq$-greatest element.
\end{proposition}
\begin{proof}
For $(\eps_R,\eps_D)$ let Radio have $(|r|\le1,\ |m-r|\le\eps_R)$ and Display have $(\mathit{true},\ |d-m|\le\eps_D)$.
Take $D_1$ with $(\eps_R,\eps_D)=(\frac15,\frac3{10})$ and $D_2$ with $(\frac3{10},\frac15)$.
Both are acceptable. (C1) is the triangle inequality: $|d-r|\le|d-m|+|m-r|\le\frac12$.
(C2) holds since Radio assumes $A_P$ and Display assumes $\mathit{true}$.
The outputs $m:=r$ and $d:=m$ show receptiveness and causality.
Suppose $D^*$ is $\sqsubseteq$-greatest. Then $S^*_i\supseteq S^1_i\cup S^2_i$ for both components.
The trace $r=0$, $m=\frac3{10}$, $d=\frac35$ lies in $A_P$, in $S^2_{\mathrm{Radio}}$ and in $S^1_{\mathrm{Display}}$.
So it lies in $A_P\wedge S^*_{\mathrm{Radio}}\wedge S^*_{\mathrm{Display}}$, but $|d-r|=\frac35>\frac12$.
Hence $D^*$ violates (C1).
\end{proof}
Z3 checks every claim of the proof; for the componentwise union it returns the mirror trace $r=0$, $m=-\frac3{10}$, $d=-\frac35$.
A choice rule is therefore needed. Ours maximises the smallest slack (Section~\ref{sec:eps}).

\section{The method for a DAG of delay-0 components}\label{sec:method}
\paragraph{Support.}
The support $N$ is the set of components with a directed path, possibly of length 0, to the owner of a variable of $G_P$.
Components outside $N$ get $(\mathit{true},\mathit{true})$.
\begin{lemma}\label{lem:outside}
Assume $A_P$ mentions only environment variables and the components outside $N$ form an acyclic graph.
Let $D$ be a valid delegation whose components outside $N$ are receptive.
Then $D'$, which gives every component outside $N$ the contract $(\mathit{true},\mathit{true})$, is valid.
\end{lemma}
\begin{proof}
No component of $N$ reads an output of a component $j\notin N$, since such an edge would put $j$ in $N$.
No variable of $G_P$ is owned by such a $j$, for the same reason.
Let $x\in A_P\wedge\bigwedge_{i\in N}S_i$. Change the outputs of the components outside $N$ in a topological order, each time choosing outputs in $S_j$ for the current inputs.
This is possible by receptiveness. The new trace $x'$ agrees with $x$ on the environment, on $\cols(\ports(i))$ for $i\in N$ and on the columns of $G_P$.
So $x'\in A_P\wedge\bigwedge_jS_j$, and (C1) for $D$ gives $x'\in G_P$, hence $x\in G_P$.
(C2) for $D'$ follows the same way, since each $A_i$ with $i\in N$ is over columns where $x$ and $x'$ agree.
\end{proof}
\paragraph{Order and column sets.}
Order $N$ topologically as $i_1,\dots,i_p$; ties are broken by name.
For $k=1,\dots,p$ let $\rest_k:=\{i_k,\dots,i_p\}$, $F_k:=\cols(\ports(i_k))$ and $K_k:=\cols(\ports(\rest_k))$.
For $k<p$ define the private columns $Z_k:=F_k\setminus K_{k+1}$, the edge variables $E_k:=\{v\in\out(i_k): v\text{ is read in }\rest_{k+1}\}$ with $e_k:=|E_k|$,
the columns $O_k:=\cols(\out(\rest_{k+1}))$, and $\env_k$, the environment variables read in $\rest_{k+1}$ but not by $i_k$.
$Z_k(t)$ and $z_k(t)$ denote the private columns at time $t$ and their values.

\paragraph{Recursion.}
Fix $\eps\ge0$. Put $Q_1:=G_P$. For $k=1,\dots,p-1$:
\begin{enumerate}[label=(\alph*)]
\item \emph{Needed forms.} Each atom of $Q_k$ at time $t$ has a coefficient row on $Z_k(t)$.
Keep the nonzero rows. Take a maximal linearly independent subfamily in order of first occurrence (Gaussian elimination over $\Q$).
Scale each row so that its first nonzero entry is 1; this convention fixes the sign and scale of the forwarded value.
The rows form $L_k(t)$, of rank $\rho_k(t)$. They are the query's own linear forms of what only $i_k$ sees.
\item \emph{Capacity.} If $\rho_k(t)>e_k$ for some $t$, stop. Report that $i_k$ must forward $\rho_k(t)$ independent forms at time $t$ but has $e_k$ edge variables, and list the forms.
Otherwise assign the $j$-th form to the $j$-th edge variable $m_{k,j}$ in name order (a labelling convention).
\item \emph{Interface.} $I_k(\eps):=\bigwedge_t\bigwedge_{j\le\rho_k(t)}|m_{k,j}(t)-L_k(t)_j\,z_k(t)|\le\eps$. Other edge columns are free.
$I_k$ is a conjunction of time-stamped atoms over $F_k$.
\item \emph{Shadow.} $\Sigma_k:=\Pi_{F_k}\,\pi_{\overline{O_k}}(A_P\to Q_k)$, where $\overline{O_k}$ is the complement of $O_k$.
First the outputs of later components are eliminated existentially, then everything else that $i_k$ does not see, universally.
\item \emph{Emit} $G_{i_k}:=I_k(\eps)\wedge\Sigma_k$ and $A_{i_k}:=\pi_{\cols(\inp(i_k))}A_P$.
\item \emph{Remaining query.} $Q_{k+1}:=\Pi_{K_{k+1}}\big((A_P\wedge I_k(\eps)\wedge\Sigma_k)\to Q_k\big)$.
\end{enumerate}
Finally $G_{i_p}:=Q_p$ and $A_{i_p}:=\pi_{\cols(\inp(i_p))}A_P$.
By Fact~\ref{fact:closure}, $\Sigma_k$, $A_{i_k}$ and $Q_{k+1}$ are PolySTL over $F_k$, $\cols(\inp(i_k))$ and $K_{k+1}$.
If $G_P$ mentions a column that no component of $N$ sees, the script uses $G_{i_p}:=\Pi_{F_p}(A_P\to Q_p)$ instead; Theorem~\ref{thm:sound} covers both.
For printing, the script writes $Q_{k+1}$ as $\Omega\to Q'$ with $\Omega:=\pi_{K_{k+1}}(A_P\wedge I_k\wedge\Sigma_k)$, projects each conjunct of $Q_k$ separately ($\forall$ distributes over $\wedge$), simplifies $Q'$ under $\Omega$, and checks the result equivalent to the definition.

The interface relations $I_k$ are the only objects of assumed form; they are derived from the query's own linear forms (across a delay, from its own atoms), and forwarding is sufficient, not necessary.
Every other obligation is whatever projection returns; no template is assumed.

\begin{theorem}[soundness]\label{thm:sound}
For every $\eps\ge0$, and for any interface formulas $I_k$ over $F_k$ in place of $I_k(\eps)$, the emitted delegation satisfies (C1) and (C2).
\end{theorem}
\begin{proof}
(C2): $A_P\models\pi_{\cols(\inp(i))}A_P=A_i$ for $i\in N$, and $A_i=\mathit{true}$ otherwise.
(C1): let $x\in A_P\wedge\bigwedge_iS_i$. By (C2), $x\in A_i$, so $x\in G_i$ for every $i$.
We show $x\in Q_k$ for $k=p,\dots,1$.
For $k=p$: $G_{i_p}=Q_p$, or $G_{i_p}=\Pi_{F_p}(A_P\to Q_p)\subseteq A_P\to Q_p$ and $x\in A_P$.
For $k<p$: $x\in Q_{k+1}\subseteq(A_P\wedge I_k\wedge\Sigma_k)\to Q_k$ because $\Pi_K\varphi\subseteq\varphi$.
Also $x\in A_P$ and $x\in G_{i_k}=I_k\wedge\Sigma_k$, so $x\in Q_k$.
Hence $x\in Q_1=G_P$. The proof never uses the form of $I_k$.
\end{proof}
Receptiveness is not guaranteed by construction. The script decides it per component by the sentence $\forall u\,\exists w.\,S_i$, and causality by the alternation of the definition.

\begin{theorem}[exact forwarding]\label{thm:exact}
Let $\eps=0$, all delays $0$ and $p\ge2$, and suppose the method does not stop.
Assume (H1) every variable of $G_P$ is a port of a component of $N$, and (H2) $A_P$ mentions only environment variables.
Say $\Sigma_k$ is \emph{nonempty on inputs} if every value $u$ of $\cols(\inp(i_k))$ with $u\in A_{i_k}$ has some $w$ with $(u,w)\in I_k(0)\wedge\Sigma_k$. Then:
\begin{enumerate}[label=(\roman*)]
\item the last component $i_p$ is receptive;
\item the delegation is receptive iff $\Sigma_k$ is nonempty on inputs for every $k<p$;
\item if the delegation is receptive, the query is centrally realizable on the support: for every value of the environment columns there are values of $\cols(\out(N))$ with $A_P\to G_P$.
If every component is causal, the central strategy can be taken causal.
\end{enumerate}
Hence the delegation is receptive iff the query is centrally realizable on the support and every $\Sigma_k$ is nonempty on inputs.
\end{theorem}
\begin{proof}
(i) Let $j:=i_p$, $k:=p-1$ and $\Phi:=A_P\wedge I_k(0)\wedge\Sigma_k$.
$\Phi$ does not mention $\cols(\out(j))$: $A_P$ by (H2), and $I_k,\Sigma_k$ are over $F_k$, which misses $\cols(\out(j))$ because $i_k$ precedes $j$.
Fix a value $u$ of $\cols(\inp(j))$.
If no trace extending $u$ lies in $\Phi$, every $w$ gives $(u,w)\in Q_p$.
Otherwise pick a trace $x_0$ extending $u$ with $x_0\in\Phi$.
From $x_0\in\Sigma_k\subseteq\pi_{\overline{O_k}}(A_P\to Q_k)$ and $O_k=\cols(\out(j))$, some trace $z$ that differs from $x_0$ only on $\cols(\out(j))$ satisfies $A_P\to Q_k$.
By (H2), $z\in A_P$, so $z\in Q_k$. Let $w:=z|_{\cols(\out(j))}$.
Let $x$ extend $(u,w)$ with $x\in\Phi$. We show $x\in Q_k$.
The columns of $Q_k$ lie in $Z_k\cup F_p$: for $k\ge2$ by construction, for $k=1$ by (H1).
The traces $x$ and $z$ agree on $F_p$. Both satisfy $I_k(0)$: $x$ by choice, $z$ because it agrees with $x_0$ on $F_k$.
The edge columns lie in $F_p$, so $L_k(t)\,x(Z_k(t))=L_k(t)\,z(Z_k(t))$ for every $t$.
The row of any atom of $Q_k$ on $Z_k(t)$ is a combination $\lambda^\top L_k(t)$, by maximality in step (a).
So every atom of $Q_k$ has the same value at $x$ and at $z$. As $z\in Q_k$, also $x\in Q_k$.
Thus $(u,w)\in Q_p$.
(ii) If $i_k$ is receptive, inputs in $A_{i_k}$ need outputs in $G_{i_k}=I_k(0)\wedge\Sigma_k$. Conversely, nonemptiness gives receptiveness of $i_k$ for $k<p$, since inputs outside $A_{i_k}$ accept any output; $i_p$ is receptive by (i), and components outside $N$ trivially.
(iii) Fix environment values satisfying $A_P$. Choose outputs of $i_1,\dots,i_p$ in turn, each in $S_{i_k}$ for the inputs already fixed.
The resulting trace lies in $A_P\wedge\bigwedge_iS_i$, so in $G_P$ by Theorem~\ref{thm:sound}.
The columns of $G_P$ lie in the environment and $\cols(\out(N))$, so the other columns do not matter.
With causal strategies, do the same at each time $t$ in the order $i_1,\dots,i_p$.
The final claim combines (ii) and (iii).
\end{proof}
\begin{conjecture}\label{conj:causal}
Under the hypotheses of Theorem~\ref{thm:exact} with $H>0$, if each $I_k(0)\wedge\Sigma_k$ has a causal strategy on $A_{i_k}$, then the last component is causal.
\end{conjecture}
The proof of (i) builds $w$ from the whole input trace, so it does not give causality; E6 with delays $(0,0)$ agrees with the conjecture.
If $p=1$ the method emits $G_P$ itself.

\paragraph{Choice of $\eps$.}\label{sec:eps}
Let $S$ be the implementation set of the contract emitted for component $i$, and $\delta\ge0$. Let $\eta$ range over values of $\cols(\out(i))$ and
$\Psi(\delta):=\mathcal{Q}_0\cdots\mathcal{Q}_H\,\forall\eta\,(|\eta|_\infty\le\delta\to S[w:=w+\eta])$, with the alternation of the causality definition.
The slack is $\mathrm{slack}(S):=\sup\{\delta\ge0:\Psi(\delta)\}$. For rational $\delta$, $\Psi(\delta)$ is a closed sentence of linear real arithmetic; Z3 decides it by \texttt{qe2}.
Let $\Theta(\eps)$ be the verdict ``every obligation of $N$ emitted at tolerance $\eps$ satisfies $\Psi(\eps)$''.
The rule is $\eps^*:=\sup\{\eps\ge0:\Theta(\eps)\}$. One $\eps$ is shared by every interface, so the rule equalises the implementation tolerance along the chain.
It assumes no shape. If $\eps^*=0$ the delegation is exact forwarding.
If $\Sigma_k=\mathit{true}$ and $L_k$ involves only inputs, the interface has slack exactly $\eps$: the output $m:=L_kz_k$ tolerates every $|\eta|\le\eps$, and no larger perturbation is tolerated.
\emph{Bisection.} The script checks $\Theta(0)$, sets $\mathit{hi}:=1$ (doubling while $\Theta(\mathit{hi})$ holds) and halves $[\mathit{lo},\mathit{hi}]$ until $\mathit{hi}-\mathit{lo}\le\frac1{1024}$.
It reports both bracketing verdicts. It then reports, as $\eps^*$, the rational with the smallest denominator in $[\mathit{lo},\mathit{hi})$, after checking that its verdict holds.
The bisection presumes that $\Theta$ is monotone on the bracket. We do not prove this in general. In E1--E4 a hand bound shows $\Theta(\eps)$ fails above the reported value, so the reported value is exact.

\section{Worked examples}\label{sec:examples}
\paragraph{E1 (relay).}
$N=\{\mathrm{Radio},\mathrm{Display}\}$, order Radio, Display, $p=2$.
For $k=1$: $F_1=\{r,m\}$, $K_2=\{m,d\}$, $Z_1=\{r\}$, $E_1=\{m\}$, $O_1=\{d\}$, $\env_1=\emptyset$.
$Q_1=G_P=(d-r\le\frac12)\wedge(-d+r\le\frac12)$.
(a) The rows on $r$ are $-1$ and $1$. After scaling, $L_1=[1]$: the form is $r$, with $\rho_1=1$.
(b) $e_1=1$, so $m$ carries $r$.
(c) $I_1(\eps)=|m-r|\le\eps$: a tolerance $\eps$ around the forwarded form.
(d) $\pi_{\overline{\{d\}}}(|r|\le1\to|d-r|\le\frac12)=\mathit{true}$ (take $d:=r$), so $\Sigma_1=\mathit{true}$.
(e) $G_{\mathrm{Radio}}=|m-r|\le\eps$ and $A_{\mathrm{Radio}}=\pi_{\{r\}}A_P=|r|\le1$.
(f) $Q_2=\Pi_{\{m,d\}}(|r|\le1\wedge|m-r|\le\eps\to|d-r|\le\frac12)$. By hand: the antecedent allows some $r$ iff $|m|\le1+\eps$.
Then the allowed $r$ fill $[\max(-1,m-\eps),\min(1,m+\eps)]$. All of them satisfy $G_P$ iff $d-\frac12\le\max(-1,m-\eps)$ and $\min(1,m+\eps)\le d+\frac12$. Hence
\[Q_2(\eps)=|m|\le1+\eps\ \to\ \big((d\le-\tfrac12\vee d-m\le\tfrac12-\eps)\wedge(d\ge\tfrac12\vee d-m\ge-\tfrac12+\eps)\big).\]
For $0\le\eps\le\frac12$ this equals the four-piece formula
$\{|m|\le1-\eps\wedge|d-m|\le\frac12-\eps\}\vee\{1-\eps<m\le1+\eps\wedge\frac12\le d\le m+\frac12-\eps\}\vee\{-1-\eps\le m<-1+\eps\wedge m-\frac12+\eps\le d\le-\frac12\}\vee\{|m|>1+\eps\}$.
Z3 confirms the equivalence with $\eps$ symbolic. $A_{\mathrm{Display}}=\pi_{\{m\}}A_P=\mathit{true}$.
\emph{Slack.} If $|m|\le1-\eps$ the admissible $d$ form $[m-\frac12+\eps,m+\frac12-\eps]$, of half-width $\frac12-\eps$.
If $1-\eps<m\le1+\eps$ they form $[\frac12,m+\frac12-\eps]$, which is wider; $m<-1+\eps$ is symmetric, and $|m|>1+\eps$ allows every $d$.
So $\mathrm{slack}(Q_2)=\frac12-\eps$, the interface has slack $\eps$, and $\eps^*=\frac14$.
The bisection prints $0$:yes, $1$:no, $\frac12$:no, $\frac14$:yes, $\frac38$:no, \dots, $\frac{257}{1024}$:no, the bracket $[\frac14,\frac{257}{1024})$ and $\eps^*=\frac14$.
At $\eps=\frac14$ the script prints exactly the hand formula:
\[G_{\mathrm{Display}}=Q_2=|m| \le \tfrac{5}{4} \rightarrow ((d \le -\tfrac{1}{2} \vee d - m \le \tfrac{1}{4}) \wedge (d \ge \tfrac{1}{2} \vee d - m \ge -\tfrac{1}{4})).\]
Emitted: Radio $(|r|\le1,\ |m-r|\le\frac14)$ and Display $(\mathit{true},\ Q_2)$.
The Display obligation is not a fixed tolerance around $m$: for $m$ near $1$ it allows any $d\in[\frac12,m+\frac14]$.
Z3 reports (C1), (C2), receptive and causal for both components.

\paragraph{E2 (fan-in).}
The environment owns $s_1,s_2$. A reads $s_1$, owns $a$; B reads $s_2$, owns $b$; Fuser reads $a,b$, owns $f$; Logger reads $f$, owns $l$. $H=0$, $A_P=\mathit{true}$, $G_P=|f-s_1-s_2|\le1$.
$N=\{A,B,\mathrm{Fuser}\}$; Logger gets $(\mathit{true},\mathit{true})$. Order A, B, Fuser.
For A: $Z_1=\{s_1\}$, $\env_1=\{s_2\}$, form $s_1$, $I_1=|a-s_1|\le\eps$, $\Sigma_1=\mathit{true}$, and $Q_2=|f-a-s_2|\le1-\eps$.
For B: form $s_2$, $I_2=|b-s_2|\le\eps$, and $Q_3=|f-a-b|\le1-2\eps$.
The Fuser's slack is $1-2\eps$, so $\eps^*=\frac13$. The bracket is $[\frac{341}{1024},\frac{171}{512})$ and the script reports $\eps^*=\frac13$.
It prints $Q_2=|a-f+s_2|\le\frac23$ and $G_{\mathrm{Fuser}}=|a+b-f|\le\frac13$.

\paragraph{E3 (three-stage chain).}
C1 reads $r$, owns $m_1$; C2 reads $m_1$, owns $m_2$; C3 reads $m_2$, owns $d$. $A_P$ and $G_P$ are those of E1.
$Q_2$ is E1's $Q_2$ with $m_1$ for $m$. Its atoms have rows on $Z_2=\{m_1\}$, so $L_2=[1]$ and $I_2=|m_1-m_2|\le\eps$: a forwarded form is forwarded again.
If $|m_2|\le1-2\eps$ the admissible $d$ form $[m_2-\frac12+2\eps,m_2+\frac12-2\eps]$, so $\mathrm{slack}(Q_3)\le\frac12-2\eps$ and $\eps^*\le\frac16$.
The bracket is $[\frac{85}{512},\frac{171}{1024})$, $\eps^*=\frac16$, and the script prints
\[G_{C3}=|m_{2}| \le \tfrac{4}{3} \rightarrow ((d \le -\tfrac{1}{2} \vee d - m_{2} \le \tfrac{1}{6}) \wedge (d \ge \tfrac{1}{2} \vee d - m_{2} \ge -\tfrac{1}{6})).\]

\paragraph{E4 (a disjunctive query).}
Relay graph, $H=0$, $A_P=|r|\le1$ and $G_P=((r\ge0\to d\ge1)\wedge(r<0\to d\le-1))\vee|d-2r|\le\frac14$.
The atoms $r\ge0$ and $r<0$ and the two atoms of $|d-2r|\le\frac14$ have rows $-1,1,-2,2$ on $Z_1=\{r\}$. They are dependent, so $\rho_1=1$, the form is $r$ and $I_1=|m-r|\le\eps$.
$\Sigma_1=\mathit{true}$. At $\eps^*=\frac1{12}$ the computed $Q_2$, exactly as the script prints it, is
\begin{multline*}
G_{\mathrm{Display}}=Q_2=|m| \le \tfrac{13}{12} \rightarrow ((d \le -1 \vee d - 2m \le \tfrac{1}{12} \vee m \ge \tfrac{1}{12}) \\ \wedge (d \le -1 \vee d - 2m \ge -\tfrac{1}{12} \vee m \ge -\tfrac{1}{12})
\wedge (d \le \tfrac{1}{4} \vee d \ge 1 \vee d - 2m \le \tfrac{1}{12}) \\ \wedge (m \ge -\tfrac{1}{12} \rightarrow (d \ge 1 \vee d - 2m \ge -\tfrac{1}{12}))).
\end{multline*}
\emph{Hand bound.} Let $m=0$ and $\eps\le\frac18$, so $r$ ranges over $[-\eps,\eps]$.
If $d\ge1$, some $r<0$ needs $|d-2r|\le\frac14$, which is impossible. If $d\le-1$, some $r\ge0$ fails the same way.
So every $r$ needs $|d-2r|\le\frac14$, that is $d\in[2\eps-\frac14,\frac14-2\eps]$.
Hence $\mathrm{slack}(Q_2)\le\frac14-2\eps$ and $\eps^*\le\frac1{12}$.
The bracket is $[\frac{85}{1024},\frac{43}{512})$ and $\Theta(\frac1{12})$ holds, so $\eps^*=\frac1{12}$.
Z3 reports (C1), (C2), receptive and causal.

\paragraph{E5 (capacity).}
The environment owns $r_1,r_2$. Radio reads both and owns $m$. Display reads $m$ and owns $d_1,d_2$. $H=0$,
$A_P=|r_1|\le1\wedge|r_2|\le1$ and $G_P=|d_1-r_1|\le\frac12\wedge|d_2-r_2|\le\frac12$.
The forms are $r_1$ and $r_2$, so $\rho_1=2>e_1=1$. The method stops with the report printed in Section~\ref{sec:repro}.
This is an obstruction for the forwarding method. It is not a proof that no acceptable delegation exists.

\section{Delays}\label{sec:delays}
With $d_{i_k}=1$, an output at $t+1$ may depend on inputs at $t$. Exact forwarding would read $m(t+1)=L\,z(t)$, a cross-time relation.
\begin{lemma}[expressiveness limit]\label{lem:limit}
Let $L\neq0$ be a row over the columns $z(t)$ and $\eps\ge0$. The set $\{|m(t+1)-L\,z(t)|\le\eps\}$ is not a PolySTL formula over $z(t)\cup\{m(t+1)\}$.
\end{lemma}
\begin{proof}
An atom over these columns is over $z(t)$ alone or over $m(t+1)$ alone. In disjunctive normal form, a PolySTL formula over them is a finite union of products $\mathcal{B}\times\mathcal{B}'$, with $\mathcal{B}$ a polyhedron of $z(t)$ values and $\mathcal{B}'$ an interval of $m(t+1)$ values.
If $\mathcal{B}\times\mathcal{B}'$ lies in the set and $\mathcal{B}'\neq\emptyset$, then $L(\mathcal{B})$ lies within $\eps$ of one point, so it is bounded.
A finite union of such products bounds $L\,z(t)$ on the set. But the set contains $z$ with $L\,z$ arbitrary, taking $m:=Lz$.
\end{proof}
So continuous information cannot cross a delayed edge in this language. Finite information can: the truth values of atoms.
\paragraph{Region forwarding.}
Let $d_{i_k}=1$. If some atom of $Q_k$ mixes $Z_k$ with other columns, the method stops and reports that forwarding that form across the delay-1 edge is not expressible in PolySTL.
Otherwise let $\mu_1,\dots,\mu_q$ be the distinct atoms of $Q_k$ on $Z_k(t)$, up to negation, in order of first occurrence and with the polarity of the first occurrence.
Let $m$ be the first edge variable in name order. Use the regions $R_c:=[2c,2c+1]$, $c=\sum_j c_j2^{j-1}$, one for each truth assignment $(c_1,\dots,c_q)$. Put
\[I_k:=\textstyle\bigwedge_{t<H}\bigwedge_{c=0}^{2^q-1}\big((\bigwedge_j\mu_j(t)^{c_j})\to m(t+1)\in R_c\big),\]
where $\mu^1:=\mu$ and $\mu^0:=\neg\mu$. Atoms at time $H$ have no later edge column and stay private.
Steps (d)--(f) are unchanged. Regions of impossible assignments are never used and are harmless.
\begin{lemma}\label{lem:regionrec}
$I_k$ is receptive and causal for $d_{i_k}=1$.
\end{lemma}
\begin{proof}
At each $t<H$ exactly one assignment $c$ holds. The output $m(t+1):=2c$ depends only on inputs at time $t$ and satisfies every conjunct for $t$, since the others have false antecedents. $m(0)$ is free.
\end{proof}
\begin{theorem}[soundness with delays]\label{thm:delay}
With region forwarding at delay-1 components, the emitted delegation satisfies (C1) and (C2), and every obligation is PolySTL over its component's ports.
\end{theorem}
\begin{proof}
The proof of Theorem~\ref{thm:sound} never uses the form of $I_k$. $I_k$ is over $F_k$: its atoms are over $Z_k(t)\subseteq F_k$ and over $m(t+1)$ with $m\in\out(i_k)$.
The other obligations are PolySTL by Fact~\ref{fact:closure}.
\end{proof}

\paragraph{E6 (relay with a deadline).}
For a deadline $k$ let $G_P(k):=\Al_{[0,2]}(|r|>\frac12\to\Ev_{[0,k]}|d|\ge\frac25)\wedge\Al_{[0,2]}(\Al_{[0,k]}|r|\le\frac12\to\Ev_{[k,k]}|d|\le\frac15)$, with $A_P=\mathit{true}$ and $H=4$ on the relay graph.
The specification's E6 is $k=2$. A trigger must get a response within $k$ steps. Without a trigger in the last $k$ steps the output must be quiet.
\begin{lemma}\label{lem:quiet}
For the first conjunct alone, Radio $(\mathit{true},\mathit{true})$ and Display $(\mathit{true},\Al_{[0,4]}d=1)$ is acceptable. It violates (C1) for the full $G_P(2)$.
\end{lemma}
\begin{proof}
With $d\equiv1$ every $\Ev_{[t,t+2]}|d|\ge\frac25$ holds. Both contracts are trivially receptive and causal.
With $r\equiv0$ the quiet conjunct needs $|d(2)|\le\frac15$.
\end{proof}
This is why the second conjunct is there: it forces the Display to use information about $r$.

\emph{Delays $(1,1)$, $k=2$.} Every atom of $Q_1$ with an $r$-column is over $r$ alone, so region forwarding applies.
At $t=0,\dots,3$ the atoms are $\mu_1=r(t)>\frac12$ and $\mu_2=r(t)<-\frac12$, so $q=2$.
The regions are $R_0=[0,1]$ (quiet), $R_1=[2,3]$ ($r>\frac12$), $R_2=[4,5]$ ($r<-\frac12$) and $R_3=[6,7]$ (impossible).
The script prints $I_{\mathrm{Radio}}$ in an equivalent simplified form, where the clause of $R_3$ is true and drops out:
$\bigwedge_{t=0}^{3}\big(r(t)<-\frac12\to m(t+1)\in R_2\big)\wedge\big(r(t)>\frac12\to m(t+1)\in R_1\big)\wedge\big(|r(t)|\le\frac12\to m(t+1)\in R_0\big)$.
$\Sigma_{\mathrm{Radio}}=\mathit{true}$. The computed Display obligation is $\Omega\to Q'$ with $\Omega=\bigwedge_{t=1}^{4}m(t)\in R_0\cup R_1\cup R_2$ and
\begin{multline*}
Q'=\textstyle\bigwedge_{t=0}^{2}\big(m(t+1)\ge2\to\Ev_{[t,t+2]}|d|\ge\frac25\big)\\
\textstyle\wedge\bigwedge_{t=0}^{1}\big(\Al_{[t+1,t+3]}m\le1\to|d(t+2)|\le\frac15\big)\wedge\big(\Al_{[3,4]}m\le1\to|d(4)|\le\frac15\big).
\end{multline*}
Inside $\Omega$, $m(t+1)\ge2$ says that $r(t)$ triggered and $m(t+1)\le1$ says that it was quiet.
The last clause has no $m(5)$: $r(4)$ is never forwarded, and the universal projection takes the worst case, $r(4)$ quiet.
The verbatim formulas are in Section~\ref{sec:repro}. Z3 reports (C1), (C2), receptive and causal for both components.

\emph{Delays $(1,1)$, $k=1$.} The same steps give the obligations of Section~\ref{sec:repro}.
Z3 finds the Display obligation receptive but not causal, and the script reports it.
The report is correct for the query itself. In any composition respecting the delays, $d(t)$ depends on $r$ only up to time $t-2$.
Fix a quiet $r(0)$. Then $d(1)$ and $d(2)$ do not depend on $r(1)$ or $r(2)$.
The quiet case $r(1)=r(2)=0$ forces $|d(1)|\le\frac15$ and $|d(2)|\le\frac15$, so these bounds hold for every $r(1)$.
A trigger at $r(1)$ then gets no response in $[1,2]$.

\emph{Delays $(0,0)$, $k=2$.} Section~\ref{sec:method} applies with the form $r(t)$ at each $t$ and $I=\bigwedge_t|m(t)-r(t)|\le\eps$.
The bisection prints $0$:yes and \emph{no} at $1,\frac12,\dots,\frac1{1024}$, so $\eps^*=0$: exact forwarding.
By hand, every $\eps>0$ fails. Take $m(0)=m(1)=0$ and $m(2)=m(3)=m(4)=\frac12$.
The value $r(2)=\frac12+\eps$ is consistent and triggers, so some $|d(t)|\ge\frac25$ with $t\in\{2,3,4\}$.
The values $r(2)=r(3)=r(4)=\frac12$ are consistent and quiet, so $|d(2)|,|d(3)|,|d(4)|\le\frac15$.
$Q_{\mathrm{Display}}$ demands both, so this input has no output. At $\eps=0$ Z3 reports (C1), (C2), receptive and causal.

\section{Which components get contracts, and what is reported}\label{sec:reports}
Components outside the support $N$ get $(\mathit{true},\mathit{true})$ (Lemma~\ref{lem:outside}). In E2 this is Logger.
Every component of $N$ gets an interface obligation $I_k\wedge\Sigma_k$, except the last, which gets the remaining query $Q_p$.
The method reports instead of emitting in five cases.
\emph{Capacity}: $\rho_k(t)>e_k$; the report names the component, the time and the forms (E5).
\emph{Inexpressible}: at a delay-1 component an atom mixes private and later columns.
\emph{No edge variable}: at a delay-1 component private atoms exist but $E_k=\emptyset$.
\emph{Non-receptive} or \emph{non-causal}: a per-component Z3 verdict on the emitted obligation (E6, $k=1$).
\emph{No tolerance}: $\Theta(0)$ fails.
The script also records any failed \texttt{qe2} cross-check and any cross-time atom in a projection; by Fact~\ref{fact:closure} the second cannot occur.

\section{Limitations}\label{sec:limits}
Forwarding is sufficient, not necessary: a coarser encoding might exist where the method stops (E5).
Receptiveness is checked, not guaranteed, for $\eps>0$. Theorem~\ref{thm:exact} needs $\eps=0$, (H1) and (H2), and causality for $H>0$ is Conjecture~\ref{conj:causal}.
The bisection presumes that $\Theta$ is monotone in $\eps$, and the scaling convention of step (a) fixes the unit of $\eps$.
Quantifier elimination is exponential in the worst case.
Cycles are not treated; a delay-1 component on every cycle would allow time expansion, which is not done.
The unbounded horizon is not treated.

\section{Reproduction}\label{sec:repro}
From the directory root run \texttt{.venv/bin/python scripts/decompose.py}. It prints every intermediate formula, runs 259 projections, each cross-checked, and exits with status 0 iff all verdicts are as expected.
In the output, \texttt{@t:} marks an atom at time $t$. Below are the printed contracts and verdicts, verbatim; the intermediate formulas are omitted.
{\fontsize{7}{8}\selectfont
\begin{verbatim}
Decomposition of a global PolySTL contract over an architecture graph
'@t:' marks an atom at time t (omitted when H = 0); |e| <= c means e <= c and -e <= c.
Lemma (relay, H = 0): universal projections of A_P -> G_P on single components
  [ok] Pi_ports(Display)(A_P -> G_P) and A_P is unsatisfiable: true
  [ok] Pi_ports(Radio)(A_P -> G_P) and A_P is unsatisfiable: true
No greatest acceptable delegation (relay): (eps_R, eps_D) = (1/5, 3/10) and (3/10, 1/5)
  delegation (eps_R, eps_D) = (1/5, 3/10):
  emitted contracts (A_i, G_i):
    Radio:  A = |r| <= 1
            G = |m - r| <= 1/5
    Display:  A = true
              G = |d - m| <= 3/10
  [ok] (C1), (C2): (true, true)
  [ok] Display receptive, causal (delay 0): (true, true)
  [ok] Radio receptive, causal (delay 0): (true, true)
  delegation (eps_R, eps_D) = (3/10, 1/5):
  emitted contracts (A_i, G_i):
    Radio:  A = |r| <= 1
            G = |m - r| <= 3/10
    Display:  A = true
              G = |d - m| <= 1/5
  [ok] (C1), (C2): (true, true)
  [ok] Display receptive, causal (delay 0): (true, true)
  [ok] Radio receptive, causal (delay 0): (true, true)
  [ok] (C1) for the componentwise union of the two: false
    Z3 counterexample: r = 0, m = -3/10, d = -3/5
  [ok] point r = 0, m = 3/10, d = 3/5 satisfies A_P, both unions, and violates G_P: true
E1 relay: r env; Radio r->m; Display m->d; A_P = |r| <= 1; G_P = |d - r| <= 1/2
    bracket: eps = 1/4 yes, eps = 257/1024 no; smallest-denominator rational in [1/4, 257/1024)
        with verdict yes: eps* = 1/4
  [ok] eps*: 1/4
  emitted contracts (A_i, G_i):
    Radio:  A = |r| <= 1
            G = |m - r| <= 1/4
    Display:  A = true
              G = |m| <= 5/4 -> ((d <= -1/2 or d - m <= 1/4) and (d >= 1/2 or d - m >= -1/4))
  [ok] (C1), (C2): (true, true)
  [ok] Radio receptive, causal (delay 0): (true, true)
  [ok] Display receptive, causal (delay 0): (true, true)
  [ok] E1 parametric Q_2 equals the expected formula: true
E2 fan-in: A s1->a; B s2->b; Fuser a,b->f; Logger f->l; A_P = true; G_P = |f - s1 - s2| <= 1
    bracket: eps = 341/1024 yes, eps = 171/512 no; smallest-denominator rational in [341/1024,
        171/512) with verdict yes: eps* = 1/3
  [ok] eps*: 1/3
  emitted contracts (A_i, G_i):
    A:  A = true
        G = |a - s1| <= 1/3
    B:  A = true
        G = |b - s2| <= 1/3
    Fuser:  A = true
            G = |a + b - f| <= 1/3
    Logger:  A = true
             G = true
  [ok] (C1), (C2): (true, true)
  [ok] A receptive, causal (delay 0): (true, true)
  [ok] B receptive, causal (delay 0): (true, true)
  [ok] Fuser receptive, causal (delay 0): (true, true)
E3 chain: C1 r->m1; C2 m1->m2; C3 m2->d; A_P = |r| <= 1; G_P = |d - r| <= 1/2
    bracket: eps = 85/512 yes, eps = 171/1024 no; smallest-denominator rational in [85/512,
        171/1024) with verdict yes: eps* = 1/6
  [ok] eps*: 1/6
  emitted contracts (A_i, G_i):
    C1:  A = |r| <= 1
         G = |m1 - r| <= 1/6
    C2:  A = true
         G = |m1 - m2| <= 1/6
    C3:  A = true
         G = |m2| <= 4/3 -> ((d <= -1/2 or d - m2 <= 1/6) and (d >= 1/2 or d - m2 >= -1/6))
  [ok] (C1), (C2): (true, true)
  [ok] C1 receptive, causal (delay 0): (true, true)
  [ok] C2 receptive, causal (delay 0): (true, true)
  [ok] C3 receptive, causal (delay 0): (true, true)
E4 disjunctive query on the relay: A_P = |r| <= 1; G_P = ((r >= 0 -> d >= 1) and (r < 0 -> d <=
    -1)) or |d - 2r| <= 1/4
    bracket: eps = 85/1024 yes, eps = 43/512 no; smallest-denominator rational in [85/1024,
        43/512) with verdict yes: eps* = 1/12
  emitted contracts (A_i, G_i):
    Radio:  A = |r| <= 1
            G = |m - r| <= 1/12
    Display:  A = true
              G = |m| <= 13/12 -> ((d <= -1 or d - 2*m <= 1/12 or m >= 1/12) and (d <= -1 or
                  d - 2*m >= -1/12 or m >= -1/12) and (d <= 1/4 or d >= 1 or
                  d - 2*m <= 1/12) and (m >= -1/12 -> (d >= 1 or d - 2*m >= -1/12)))
  [ok] (C1), (C2): (true, true)
  [ok] Radio receptive, causal (delay 0): (true, true)
  [ok] Display receptive, causal (delay 0): (true, true)
E5 capacity: Radio r1,r2->m; Display m->d1,d2; A_P = |r1| <= 1 and |r2| <= 1; G_P = |d1 - r1| <=
    1/2 and |d2 - r2| <= 1/2
    STOP: component Radio must forward 2 independent linear forms of {r1, r2} at time 0 but has
        only 1 edge variable (forms: r1, r2); not delegable by forwarding
  [ok] E5 stops with the forwarding obstruction: true
E6 Lemma: with only the first conjunct, Radio (true, true), Display (true, always d = 1) is
    acceptable
  emitted contracts (A_i, G_i):
    Radio:  A = true
            G = true
    Display:  A = true
              G = @0: 1 <= d <= 1 and @1: 1 <= d <= 1 and @2: 1 <= d <= 1 and
                  @3: 1 <= d <= 1 and @4: 1 <= d <= 1
  [ok] (C1), (C2): (true, true)
  [ok] Display receptive, causal (delay 1): (true, true)
  [ok] Radio receptive, causal (delay 1): (true, true)
  the same delegation against the full G_P (with the quiet conjunct):
  [ok] (C1) for the full G_P: false
E6 relay with deadline k = 2, delays (1, 1), H = 4, A_P = true; G_P = alw[0,2](|r| > 1/2 ->
    ev[0,2] |d| >= 2/5) and alw[0,2](alw[0,2] |r| <= 1/2 -> ev[2,2] |d| <= 1/5)
  emitted contracts (A_i, G_i):
    Radio:  A = true
            G = (@0: r < -1/2 -> @1: 4 <= m <= 5) and (@0: r > 1/2 -> @1: 2 <= m <= 3) and
                (@0: |r| <= 1/2 -> @1: 0 <= m <= 1) and (@1: r < -1/2 -> @2: 4 <= m <= 5) and
                (@1: r > 1/2 -> @2: 2 <= m <= 3) and (@1: |r| <= 1/2 -> @2: 0 <= m <= 1) and
                (@2: r < -1/2 -> @3: 4 <= m <= 5) and (@2: r > 1/2 -> @3: 2 <= m <= 3) and
                (@2: |r| <= 1/2 -> @3: 0 <= m <= 1) and (@3: r < -1/2 -> @4: 4 <= m <= 5) and
                (@3: r > 1/2 -> @4: 2 <= m <= 3) and (@3: |r| <= 1/2 -> @4: 0 <= m <= 1)
    Display:  A = true
              G = (@1: 0 <= m <= 5 and @2: 0 <= m <= 5 and @3: 0 <= m <= 5 and
                  @4: 0 <= m <= 5 and (@1: m <= 1 or @1: m >= 2) and (@1: m <= 3 or
                  @1: m >= 4) and (@2: m <= 1 or @2: m >= 2) and (@2: m <= 3 or
                  @2: m >= 4) and (@3: m <= 1 or @3: m >= 2) and (@3: m <= 3 or
                  @3: m >= 4) and (@4: m <= 1 or @4: m >= 2) and (@4: m <= 3 or
                  @4: m >= 4)) -> ((@1: m >= 2 -> (@0: |d| >= 2/5 or @1: |d| >= 2/5 or
                  @2: |d| >= 2/5)) and ((@1: m <= 1 and @2: m <= 1 and @3: m <= 1) ->
                  @2: |d| <= 1/5) and (@2: m >= 2 -> (@1: |d| >= 2/5 or @2: |d| >= 2/5 or
                  @3: |d| >= 2/5)) and ((@2: m <= 1 and @3: m <= 1 and @4: m <= 1) ->
                  @3: |d| <= 1/5) and (@3: m >= 2 -> (@2: |d| >= 2/5 or @3: |d| >= 2/5 or
                  @4: |d| >= 2/5)) and ((@3: m <= 1 and @4: m <= 1) -> @4: |d| <= 1/5))
  [ok] (C1), (C2): (true, true)
  [ok] Radio receptive, causal (delay 1): (true, true)
  [ok] Display receptive, causal (delay 1): (true, true)
E6 relay with deadline k = 1, delays (1, 1), H = 4, A_P = true; G_P = alw[0,2](|r| > 1/2 ->
    ev[0,1] |d| >= 2/5) and alw[0,2](alw[0,1] |r| <= 1/2 -> ev[1,1] |d| <= 1/5)
  emitted contracts (A_i, G_i):
    Radio:  same contract as for k = 2 (Z3: equivalent)
    Display:  A = true
              G = (@1: 0 <= m <= 5 and @2: 0 <= m <= 5 and @3: 0 <= m <= 5 and
                  @4: 0 <= m <= 5 and (@1: m <= 1 or @1: m >= 2) and (@1: m <= 3 or
                  @1: m >= 4) and (@2: m <= 1 or @2: m >= 2) and (@2: m <= 3 or
                  @2: m >= 4) and (@3: m <= 1 or @3: m >= 2) and (@3: m <= 3 or
                  @3: m >= 4) and (@4: m <= 1 or @4: m >= 2) and (@4: m <= 3 or
                  @4: m >= 4)) -> ((@1: m >= 2 -> (@0: |d| >= 2/5 or @1: |d| >= 2/5)) and
                  ((@1: m <= 1 and @2: m <= 1) -> @1: |d| <= 1/5) and (@2: m >= 2 ->
                  (@1: |d| >= 2/5 or @2: |d| >= 2/5)) and ((@2: m <= 1 and @3: m <= 1) ->
                  @2: |d| <= 1/5) and (@3: m >= 2 -> (@2: |d| >= 2/5 or @3: |d| >= 2/5)) and
                  ((@3: m <= 1 and @4: m <= 1) -> @3: |d| <= 1/5))
  Display receptive (non-causal): true
  [ok] Display causal (delay 1) at k = 1: false
E6 with delays (0, 0), deadline k = 2, H = 4: continuous forwarding of r
    bracket: eps = 0 yes, eps = 1/1024 no; smallest-denominator rational in [0, 1/1024) with
        verdict yes: eps* = 0
  emitted contracts (A_i, G_i):
    Radio:  A = true
            G = @0: |m - r| <= 0 and @1: |m - r| <= 0 and @2: |m - r| <= 0 and
                @3: |m - r| <= 0 and @4: |m - r| <= 0
    Display:  A = true
              G = ((@0: |m| <= 1/2 and @1: |m| <= 1/2 and @2: |m| <= 1/2) ->
                  @2: |d| <= 1/5) and (@0: |d| >= 2/5 or @0: m <= 1/2 or @1: |d| >= 2/5 or
                  @2: |d| >= 2/5) and (@0: |d| >= 2/5 or @0: m >= -1/2 or @1: |d| >= 2/5 or
                  @2: |d| >= 2/5) and ((@1: |m| <= 1/2 and @2: |m| <= 1/2 and
                  @3: |m| <= 1/2) -> @3: |d| <= 1/5) and (@1: |d| >= 2/5 or @1: m <= 1/2 or
                  @2: |d| >= 2/5 or @3: |d| >= 2/5) and (@1: |d| >= 2/5 or @1: m >= -1/2 or
                  @2: |d| >= 2/5 or @3: |d| >= 2/5) and ((@2: |m| <= 1/2 and
                  @3: |m| <= 1/2 and @4: |m| <= 1/2) -> @4: |d| <= 1/5) and
                  (@2: |d| >= 2/5 or @2: m <= 1/2 or @3: |d| >= 2/5 or @4: |d| >= 2/5) and
                  (@2: |d| >= 2/5 or @2: m >= -1/2 or @3: |d| >= 2/5 or @4: |d| >= 2/5)
  [ok] (C1), (C2): (true, true)
  [ok] Radio receptive, causal (delay 0): (true, true)
  [ok] Display receptive, causal (delay 0): (true, true)
projections computed: 259 (each cross-checked against its quantified input)
all verdicts as expected
\end{verbatim}
}

{\small\begin{thebibliography}{9}
\bibitem{contracts} A.~Benveniste, B.~Caillaud, D.~Nickovic, R.~Passerone, J.-B.~Raclet, P.~Reinkemeier, A.~Sangiovanni-Vincentelli, W.~Damm, T.~A.~Henzinger, K.~G.~Larsen.
Contracts for system design. \emph{Foundations and Trends in Electronic Design Automation} 12(2--3):124--400, 2018.
\bibitem{z3} L.~de~Moura, N.~Bj{\o}rner. Z3: an efficient SMT solver. TACAS 2008, LNCS 4963, 337--340.
\bibitem{qe2} N.~Bj{\o}rner, M.~Janota. Playing with quantified satisfaction. LPAR-20 short papers, EPiC 35, 2015.
\end{thebibliography}}
\end{document}
